% Generated draft: requires mathematical and source review.
\documentclass[11pt]{article}
\usepackage[margin=25mm]{geometry}
\usepackage{fontspec}
\setmainfont{Latin Modern Roman}
\usepackage{amsmath,amssymb,mathtools,unicode-math}
\setmathfont{Latin Modern Math}
\usepackage{enumitem,needspace,xurl,hyperref}
\setlength{\parindent}{0pt}
\setlength{\parskip}{5pt}
\setlength{\emergencystretch}{3em}
\newcommand{\N}{\mathbb{N}}
\newcommand{\Z}{\mathbb{Z}}
\newcommand{\Q}{\mathbb{Q}}
\newcommand{\R}{\mathbb{R}}
\newcommand{\C}{\mathbb{C}}
\newcommand{\F}{\mathbb{F}}
\newcommand{\A}{\mathbb{A}}
\newcommand{\PP}{\mathbb{P}}
\newcommand{\E}{\mathbb{E}}
\newcommand{\eps}{\varepsilon}
\newcommand{\dd}{\,\mathrm{d}}
\begin{document}
\section*{UnsolvedMath ID: 8800087}
\textbf{Source problem code:} catalogue entry
\subsection{Definitions and mathematical statement}
The source asks for a baby-step/giant-step infrastructure framework that speeds infrastructure discrete logarithms or point counting by a polynomial factor. In algebraic-number/function-field computation, an infrastructure is typically a collection of reduced objects with distance information and reduction-based baby and giant steps; it is not automatically a finite group with an associative operation.

For comparison, in a cyclic group of order \(N\), ordinary baby-step/giant-step solves a discrete logarithm using on the order of \(\sqrt N\) group operations and comparable stored group elements. Infrastructure analogues must account for distance corrections, reduction costs and the relevant search interval or regulator. The excerpt does not specify the number/function-field class, unit rank, infrastructure encoding, distance problem, or step-operation cost. A rank-one distance-circle model is therefore not silently imposed on every intended instance.

The target is a provable or otherwise explicitly qualified improvement over a specified baseline for one of the intended tasks. ``Polynomial factor'' needs a parameter: a factor polynomial in input bit length differs from a factor polynomial in the search-space size. Running time, memory, precomputation and heuristic versus worst-case guarantees also need definition. Point counting may refer to curve points, a Jacobian order or another class-number-related quantity; the source does not choose among these.

The summary reports square-root-style infrastructure algorithms and refinements by Scheidler--Williams--Zhang, Jauch--Jacobson and others, while saying no decisive general polynomial-factor speedup framework was verified. It supplies no exact complexities or theorem citations. This partial-progress report is retained without independently certifying a lower bound or current algorithmic status.
\subsection{Short English statement}
Improve infrastructure discrete-logarithm or point-counting algorithms through baby-step/giant-step ideas. The infrastructure class, target quantity, baseline and parameter for a polynomial-factor improvement are missing, so no exact complexity theorem is yet specified.
\subsection{Sources}
\begin{itemize}
\item Local registry: \url{data/unsolvedmath/problems.json}, permanent UnsolvedMath ID \(8800087\), source code catalogue entry; supplied statement, background, and research summary, updated 2026-08-06. Attributed to Renate Scheidler.
\item IPAM, \emph{Workshop on Number Theory and Cryptography---Open Problems} (2006), attachment 6665, PDF pages 184--188, printed slide 29: \url{https://www.ipam.ucla.edu/programs/workshops/workshop-i-number-theory-and-cryptography-open-problems/?tab=schedule}. Supplied PDF/legacy-PPT visual audit accessed 2026-07-29; it reports 19 inspected surviving attachments.
\item Scheidler--Williams--Zhang and Jauch--Jacobson are named by the summary without full references or exact algorithmic bounds.
\end{itemize}
\end{document}
